% section: 漸化式を解く前に

\section{漸化式を解く前に}\label{sec:before}
漸化式を解く前に,根底にある発想について触れておこう.
多くの漸化式は係数などを無視して書けば$a_{n+1}=a_n$のような形をしている.
ここに係数や数式が足されることでどんどん問題は複雑化していく.
しかしながら不変の事実として,漸化式はある数列の2項の関係を示している.
だから,2行の間で"きれい"な関係が成り立っていればその漸化式は解けるのである.
この"きれい"というのは,解いていくにつれて理解されるものであるが,その多くは
\begin{itemize}
  \item 基本4パターン（特に等比型）になっている
  \item 左辺が$n+1$に関する式,右辺が$n$に関する式になっている
  \item 左辺と右辺が同じ値
\end{itemize}
などである.
最終的にこういった形を目指しながら式変形をしているということを頭の片隅に置いておくといいだろう.
この発想は応用パターンの節で頻繁に用いられることになる.

もちろん,一般項を予測し帰納法によって証明する方法も有効である.
帰納法を使って証明するかどうかに関わらず,悩んだらとりあえず$a_5$ぐらいまで書き出してみる,というのは頭の中に入れておくべきことである.
実験をするのは漸化式に限らず数学全般において重要な概念である.
